\paragraph{\texorpdfstring{$A_4$}{A4}/\texorpdfstring{$S_3$}{S3} 中间覆盖的分歧刚性、标准表示因子与定点谱}
在既定记号下，记
\[
X:=\widetilde{C},\qquad
L=\QQ(X),\qquad
\mathrm{Gal}(L/\QQ(y))\cong S_4,
\]
并令
\[
X_A:=X/A_4,\qquad E:=X/S_3.
\]
其中 \(S_3\le S_4\) 为字母稳定子，且 \(X\to \PP^1_y\) 的惯性型为
\((C_4;C_2,C_2,C_2,C_2,C_2)\)。

\begin{conclusion}[\texorpdfstring{$A_4$}{A4} 子覆盖的单点分歧刚性]\label{con:xi-terminal-zm-a4-subcover-single-branch-rigidity}
自然商映射
\[
f:X\longrightarrow X_A=X/A_4
\]
为次数 \(12\) 的 \(A_4\)-Galois 覆盖。
其分歧点在 \(X_A\) 上恰有一处；该唯一分歧点的惯性群同构于 \(C_2\)，并由双换位生成。
\end{conclusion}
\begin{proof}[证明要点]
记 \(\varphi:X\to\PP^1_y\) 为 \(S_4\)-正则覆盖。对 \(P\in X\)，中间覆盖的惯性满足
\[
I_P(f)=I_P(\varphi)\cap A_4.
\]
对五个有限分歧值，\(I_P(\varphi)\) 由换位生成，属于奇置换，故
\(
I_P(\varphi)\cap A_4=\{1\}
\)；
于是 \(f\) 在有限分歧值上方处处不分歧。

在无穷远处分歧，\(I_P(\varphi)\cong C_4\) 由四循环 \(\rho\) 生成。
其平方 \(\rho^2\) 为双换位，属于 \(A_4\)；
且 \(A_4\) 中无阶 \(4\) 元，故
\[
I_P(\varphi)\cap A_4=\langle \rho^2\rangle\cong C_2.
\]
因此 \(f\) 仅可能在 \(\infty\) 上方分歧。

又 \(X_A\to \PP^1_y\) 为符号二次子覆盖，在 \(\infty\) 处分歧，故 \(\infty\) 在 \(X_A\) 上方仅有一个点。
综上，\(f\) 的分歧点唯一，惯性为双换位生成的 \(C_2\)。证毕。
\end{proof}

\begin{conclusion}[\texorpdfstring{$S_3$}{S3} 子覆盖 \texorpdfstring{$X\to E$}{X to E} 的分歧除子全刻画]\label{con:xi-terminal-zm-s3-subcover-branch-divisor-complete}
令
\[
\pi:X\longrightarrow E=X/S_3
\]
为次数 \(6\) 的 \(S_3\)-Galois 覆盖，则：
\begin{enumerate}
  \item \(\pi\) 的所有分歧均为简单二分歧（惯性同构于 \(C_2\)），不存在阶 \(3\) 惯性；
  \item \(\pi\) 在 \(E\) 上恰有 \(10\) 个分歧点；
  \item 对每个有限分歧值 \(y_0\in\PP^1_y\)，纤维 \(E_{y_0}\) 在 \(E\to\PP^1_y\) 中呈 \(2+1+1\) 型：其中唯一二分歧点上方 \(\pi\) 不分歧，而两个不分歧点上方 \(\pi\) 恰好分歧；
  \item 在 \(E\) 的无穷远分歧点上方，\(\pi\) 不分歧。
\end{enumerate}
\end{conclusion}
\begin{proof}[证明要点]
设 \(H=S_3\le G=S_4\)。则
\[
\QQ(y)\subset \QQ(E)=L^H\subset L=\QQ(X),
\]
且 \(\mathrm{Gal}(L/L^H)=H\)。

取有限分歧值 \(y_0\)。在 \(\varphi\) 中其惯性为换位 \(\tau\in S_4\)。
考察 \(\tau\) 在陪集空间 \(S_4/H\simeq\{1,2,3,4\}\) 上的作用，循环分解为 \(2+1+1\)。
这对应 \(E\to\PP^1_y\) 在 \(y_0\) 上有一个二分歧点和两个不分歧点。

若 \(R\in E_{y_0}\)，取 \(P\in X\) 在其上方，则
\[
I_P(\pi)=H\cap g\langle\tau\rangle g^{-1}
\]
（\(g\) 由点选择决定）。
当 \(R\) 对应 \(\tau\) 的不动点时，\(g^{-1}\tau g\in H\)，故 \(I_P(\pi)\cong C_2\)；
当 \(R\) 对应长度 \(2\) 轨道时，\(g^{-1}\tau g\notin H\)，故 \(I_P(\pi)=\{1\}\)。
于是每个有限分歧值贡献恰好两处 \(\pi\)-分歧点。

在无穷远处，\(\varphi\) 的惯性为 \(C_4=\langle\rho\rangle\)。
对任意 \(g\)，\(H\cap gC_4g^{-1}=\{1\}\)：因为 \(H\cong S_3\) 不含阶 \(4\) 元，且其阶 \(2\) 元均为换位，而 \(gC_4g^{-1}\) 的唯一阶 \(2\) 元是双换位。
故 \(\pi\) 在无穷远上方不分歧。

因此总分歧点数为 \(5\times2=10\)，且惯性均为 \(C_2\)。
最后由 Riemann--Hurwitz（\(g(E)=1\)）
\[
2g(X)-2=\sum_{P\in X}(e_P-1).
\]
每个分歧点在 \(E\) 上方有 \(6/2=3\) 个 ramification 点，每点贡献 \(1\)，故总贡献 \(10\times3=30\)，与 \(g(X)=16\) 一致，并排除额外分歧与阶 \(3\) 惯性。证毕。
\end{proof}

\begin{conclusion}[标准表示因子的几何实现与 \texorpdfstring{$A_{V_3}\sim_{\QQ}E^3$}{AV3 isogenous to E^3}]\label{con:xi-terminal-zm-standard-factor-geometric-realization-e-cube}
令 \(H_1,\dots,H_4\le S_4\) 为四个共轭字母稳定子（均同构于 \(S_3\)），并记
\[
E_i:=X/H_i\qquad(1\le i\le4).
\]
则各 \(E_i\) 皆为椭圆曲线，且在字母重标下两两同构。
记 \(\pi_i:X\to E_i\)，固定 \(O_i\in E_i\)，定义同态
\[
\Phi:E_1\times E_2\times E_3\times E_4\longrightarrow \operatorname{Jac}(X),\qquad
(P_1,\dots,P_4)\longmapsto \sum_{i=1}^4 \pi_i^*(P_i-O_i).
\]
令 \(A_{V_3}:=\operatorname{Im}(\Phi)\)。则：
\begin{enumerate}
  \item \(A_{V_3}\) 为 \(\operatorname{Jac}(X)\) 的 \(S_4\)-稳定阿贝尔子簇，其余切表示同构于标准不可约表示 \(V_3\)；
  \item \(\ker(\Phi)^0\) 为对角嵌入椭圆曲线
  \[
  E_\Delta=\{(P,P,P,P)\}\subset E_1\times\cdots\times E_4;
  \]
  \item 存在 \(S_4\)-等变同源
  \[
  (E_1\times\cdots\times E_4)/E_\Delta\ \sim_{\QQ}\ A_{V_3},
  \]
  从而（忽略标号同构 \(E_i\cong E\)）
  \[
  A_{V_3}\sim_{\QQ}E^3.
  \]
\end{enumerate}
\end{conclusion}
\begin{proof}[证明要点]
在余切空间上，\(\Phi\) 的微分为
\[
d\Phi:\bigoplus_{i=1}^4 H^0(E_i,\Omega^1)\longrightarrow H^0(X,\Omega^1).
\]
左端是四点置换表示，分解为 \(\mathbf 1\oplus V_3\)。
由 \(X/S_4\simeq \PP^1_y\) 得
\[
H^0(X,\Omega^1)^{S_4}\cong H^0(\PP^1,\Omega^1)=0,
\]
故像中无平凡分量。
另一方面，每个 \(\pi_i^*:H^0(E_i,\Omega^1)\to H^0(X,\Omega^1)\) 非零，故 \(d\Phi\) 像非零。
于是 \(\operatorname{Im}(d\Phi)\) 必为 \(V_3\) 分量，维数 \(3\)；从而 \(A_{V_3}\) 维数为 \(3\)，且余切表示为 \(V_3\)。

又 \(\dim(E_1\times\cdots\times E_4)=4\)，故 \(\ker(\Phi)^0\) 维数为 \(1\)。
其切空间为置换表示中的唯一 \(S_4\)-稳定一维子空间，即对角线 \(\mathbf 1\)；
因此 \(\ker(\Phi)^0=E_\Delta\)。

\(\Phi\) 于是分解为
\[
E_1\times\cdots\times E_4 \twoheadrightarrow (E_1\times\cdots\times E_4)/E_\Delta
\xrightarrow{\ \overline{\Phi}\ } A_{V_3},
\]
且 \(\overline{\Phi}\) 在余切空间上同构，故核有限，\(\overline{\Phi}\) 为同源。
识别 \(E_i\cong E\) 即得 \(A_{V_3}\sim_{\QQ}E^3\)。证毕。
\end{proof}

\begin{conclusion}[\texorpdfstring{$A_4$}{A4} 作用的定点谱：双换位固定两点而三循环无定点]\label{con:xi-terminal-zm-a4-fixedpoint-spectrum-v4-c3}
设 \(V_4\lhd A_4\) 为 Klein 四元子群。则：
\begin{enumerate}
  \item 对任意非平凡双换位 \(\delta\in V_4\setminus\{1\}\)，有
  \[
  \#\mathrm{Fix}_X(\delta)=2;
  \]
  \item 对任意三循环 \(\sigma\in A_4\)（阶 \(3\) 元），有
  \[
  \mathrm{Fix}_X(\sigma)=\varnothing.
  \]
\end{enumerate}
\end{conclusion}
\begin{proof}[证明要点]
由结论 \ref{con:xi-terminal-zm-a4-subcover-single-branch-rigidity}，\(f:X\to X_A\) 仅在一点 \(Q\in X_A\) 分歧，且惯性 \(I\cong C_2\)。
故任一非平凡群元若有定点，该定点必位于 \(f^{-1}(Q)\)。
又
\[
\#f^{-1}(Q)=\frac{|A_4|}{|I|}=\frac{12}{2}=6.
\]

取 \(\delta\in I\setminus\{1\}\)。其在陪集 \(A_4/I\) 上固定陪集数为
\[
\#\{gI:\delta gI=gI\}
=\frac{|N_{A_4}(I)|}{|I|}.
\]
而 \(N_{A_4}(I)=V_4\)，故固定陪集数为 \(4/2=2\)。
这正对应 \(\delta\) 在 \(f^{-1}(Q)\subset X\) 上的定点数；在其余点无定点，故该 \(\delta\) 满足 \(\#\mathrm{Fix}_X(\delta)=2\)。
又 \(A_4\) 中三个非平凡双换位彼此共轭，而定点数按共轭类不变，故结论对任意 \(\delta\in V_4\setminus\{1\}\) 成立。

若 \(\sigma\) 为三循环，则 \(\sigma\notin V_4\)，且不可能落入任何阶 \(2\) 惯性子群。
由于分歧仅来自阶 \(2\) 惯性，\(\sigma\) 不能稳定任何点，故 \(\mathrm{Fix}_X(\sigma)=\varnothing\)。证毕。
\end{proof}

\endinput
